\section{Preliminaries}

\begin{definition}[Generalized Reed--Solomon code]
    Let $q$ be a prime power, $n \leq q$, $1 \leq k \leq n$. Choose distinct
    $x_1,\dots,x_n \in \mathbb{F}_q$ (evaluation points) and nonzero
    $y_1,\dots,y_n
        \in \mathbb{F}_q^*$ (column multipliers). The \emph{generalized
        Reed--Solomon
        code} $\mathrm{GRS}_{n,k}(x,y)$ is the $[n,k]$ linear code with
    generator
    matrix
    \[
        (G_{\mathrm{GRS}})_{i,j} = y_j x_j^i, \qquad i = 0,\dots,k-1,\; j =
        1,\dots,n,
    \]
    i.e., the evaluation of all polynomials of degree $< k$ at the $x_j$,
    scaled by
    $y_j$. GRS codes are MDS: minimum distance $d = n-k+1$.
    See~\cite[Ch.~10]{MS77}
    or~\cite[Ch.~5]{HP03}.
\end{definition}

\begin{definition}[Shortening]
    Let $C$ be an $[n,k]$ code and $L \subset \{1,\dots,n\}$. The
    \emph{shortening}
    of $C$ at $L$, denoted $C_L$ or $\mathcal{S}_L(C)$, is obtained by
    restricting to
    codewords that vanish on $L$, then deleting those positions. If the
    generator
    columns at $L$ are linearly independent, $\dim C_L = k - |L|$ and the
    length is
    $n - |L|$.
\end{definition}

\subsection{The RLCE scheme}

We use the $S \cdot G_s \cdot P_1 \cdot A \cdot P_2$ formulation from the RLCE
specification~\cite{Wan16,RLCE-spec}. Everything happens over the finite field $\Fq$ for 
$q=2^m$. 

Let $G_s$ be a $k \times n$ GRS generator
matrix. Key generation:
\begin{itemize}
    \item $S$: uniform random $k \times k$ invertible matrix (scrambler).
    \item $P_1$: uniform random $n \times n$ permutation matrix.
    \item $A$: $n \times (n+w)$ \emph{expansion} matrix. For $w$ designated
          positions $j_1,\dots,j_w$, column $j_t$ of $G_s P_1$ is replaced by
          two columns
          \[
              a_t g_{j_t} + c_t r_t, \qquad b_t g_{j_t} + d_t r_t,
          \]
          where $r_t \in \mathbb{F}_q^k$ is uniform random and
          $\bigl(\begin{smallmatrix}a_t&b_t\\c_t&d_t\end{smallmatrix}\bigr) \in
              \mathrm{GL}_2(\mathbb{F}_q)$ is uniform random. All other columns
          pass through
          unchanged. The pairs $(a_t g_{j_t} + c_t r_t,\; b_t g_{j_t} + d_t
              r_t)$ are the
          \emph{twin pairs}.
    \item $P_2$: uniform random $(n+w) \times (n+w)$ permutation matrix.
\end{itemize}
The public key is $G_{\mathrm{pub}} = S \cdot G_s \cdot P_1 \cdot A \cdot P_2$,
a $k \times (n+w)$ matrix. The secret key is $(S, G_s, P_1, A, P_2)$.
Encryption is $c = m G_{\mathrm{pub}} + e$ with $\mathrm{wt}(e) \leq t$;
decryption unpermutes by $P_2^{-1}$, unmixes each twin pair via the $2 \times
    2$ inverse, unpermutes by $P_1^{-1}$, and decodes the GRS code. Security of
the post-CLT paper updated version
relies
on the indistinguishability of $G_{\mathrm{pub}}$ from random solely at $w =
    n-k$; for a more precise statement of
security, consult the specification.


For completeness, we reproduce the specification's parameter sets below.

\begin{table}[htbp]
\centering
\footnotesize
\begin{tabular}{lccccc}
\toprule
Set &  n & k & t & w & m \\
\midrule
ID-0 & 630 & 470 & 80 & 160 & 10 \\
ID-1 & 1000 & 764 & 118 & 236 & 10 \\
ID-2 & 1360 & 800 & 280 & 560 & 11 \\
\bottomrule
\end{tabular}
\caption{Parameter sets from the RLCE specification}
\label{tab:spec-params}
\end{table}
\subsection{The CLT Attack}

We work with $G_{\mathrm{GRS}}$ in canonical form as above. The RLCE public key
has $N = n + w$ columns: $n - w$ pure GRS columns and $w$ twin pairs, each
$(c_i, c_{\tau(i)})$ of the form
\[
    c_i = a G_j + c R, \qquad c_{\tau(i)} = b G_j + d R,
\]
with $G_j$ a GRS column, $R \in \mathbb{F}_q^k$ random, and
$\bigl(\begin{smallmatrix}a&b\\c&d\end{smallmatrix}\bigr)$ invertible, all
permuted by secret $\pi$.

Shortening at $L$ deletes those positions after restricting to codewords
vanishing there. By~\cite[Lemma~24]{CLT18}, shortening at one twin-pair member
\emph{derandomizes} the other. If $L$ hits every pair, $C_L = M \cdot
    G'_{\mathrm{GRS}}$ with $M$ generic $k' \times k_{\mathrm{GRS}}$, $k' = k -
    |L|$. If $L$ misses $j$ pairs, those $2j$ columns remain mixed (noisy).

The ambient GRS dimension is $k_{\mathrm{GRS}} = k' + w - j$ for $j$ unhit
pairs.

\paragraph{High-level overview.} The Couvreur--Lequesne--Tillich
attack~\cite{CLT18} proceeds in three phases. First, \emph{distinguish}: the
Schur square $C^{\star 2}$ has abnormally low dimension when the gap $c = (n-k)
    - w$ exceeds $(3+\sqrt{16w+1})/2$, revealing twin-pair positions. Second,
\emph{unmix}: for each identified pair, the $2 \times 2$ mixing is inverted to
recover pure GRS columns. Third, \emph{recover}:
Sidelnikov--Shestakov~\cite{sidelnikov1992insecurity} on the
reconstructed GRS code yields the secret support and multipliers. The attack
\emph{succeeds} for $w$ sufficiently below $n-k$ but \emph{fails} at the
shipped $w = n-k$ boundary, where the square is provably full and no
distinguisher was known.

\subsection{Hasse Derivatives, Multiplicity Codes, and Jets}\label{sec:hasse}

We follow the treatment of~\cite{GIJS26}. For $\mathbf{a} = (a_1,\dots,a_{k'})
    \in \mathbb{Z}_{\geq 0}^{k'}$ write $|\mathbf{a}| = \sum_i a_i$ for its
$\ell_1$-norm, and let
\[
    A_{<s,k'} := \{\mathbf{a} \in \mathbb{Z}_{\geq 0}^{k'} : |\mathbf{a}| < s\}
\]
be the set of multi-indices of order below $s$.  We write
$\mathbf{X}^{\mathbf{a}} := \prod_i X_i^{a_i}$ for the corresponding monomial.

\begin{definition}[Hasse derivatives~{\cite[Def.~2.7]{GIJS26}}]
    Let $F \in \mathbb{F}[X_1,\dots,X_{k'}]$ over a field $\mathbb{F}$.  The
    \emph{Hasse derivative} of $F$ with respect to
    $\mathbf{a} \in \mathbb{Z}_{\geq 0}^{k'}$ is
    \[
        \partial^{[\mathbf{a}]} F
        = \sum_{\mathbf{b} \geq \mathbf{a}}
        \binom{\mathbf{b}}{\mathbf{a}} \cdot \mathrm{Coeff}_{\mathbf{b}}(F)\,
        \mathbf{X}^{\mathbf{b}-\mathbf{a}},
    \]
    where $\mathrm{Coeff}_{\mathbf{b}}(F)$ is the coefficient of
    $\mathbf{X}^{\mathbf{b}}$ in $F$ and
    $\binom{\mathbf{b}}{\mathbf{a}} = \prod_i \binom{b_i}{a_i}$.  Equivalently,
    Hasse derivatives are the finite-field analogue of ordinary partial
    derivatives, characterized by the formal Taylor expansion
    \[
        F(\mathbf{X}+\mathbf{Z})
        = \sum_{\mathbf{a}} (\partial^{[\mathbf{a}]}
        F)(\mathbf{X})\,\mathbf{Z}^{\mathbf{a}}.
    \]
    For a univariate polynomial $f(Z) = \sum_{j=0}^{d} c_j Z^j$ we write
    $\partial^{[r]} f(Z) = \sum_{j=r}^{d} \binom{j}{r} c_j Z^{j-r}$; in
    characteristic $0$ this equals $f^{(r)}(Z)/r!$.  The Hasse-derivative chain
    rule holds in any characteristic, which is what makes multiplicities usable
    over finite fields.
\end{definition}

\begin{definition}[Hasse jet~{\cite[Def.~2.9]{GIJS26}}]
    For $\mathbf{x} \in \mathbb{F}_q^{k'}$ and $s \geq 1$, the \emph{Hasse jet
        of
        order $<s$} of $F$ at $\mathbf{x}$ (also the \emph{order-$s$
        evaluation} of
    $F$ at $\mathbf{x}$) is the vector of all Hasse derivatives of order below
    $s$ evaluated at $\mathbf{x}$:
    \[
        J_{\mathbf{x}}^{<s}(F) := \bigl((\partial^{[\mathbf{a}]}
        F)(\mathbf{x})\bigr)_{\mathbf{a} \in A_{<s,k'}}.
    \]
    We say $F$ \emph{vanishes with multiplicity} (at least) $s$ at $\mathbf{x}$
    if $J_{\mathbf{x}}^{<s}(F) = 0$.
\end{definition}

\begin{definition}[Multiplicity code~{\cite[Def.~2.10]{GIJS26},
    \cite{DBLP:journals/jacm/KoppartySY14}}]
    The multiplicity code of order-$s$ evaluations of polynomials of total
    degree
    at most $d$ in $k'$ variables over $\mathbb{F}_q$ has length $q^{k'}$ over
    the
    alphabet $\Sigma = \mathbb{F}_q^{\binom{k'+s-1}{s-1}}$, indexed by
    $\mathbf{x} \in \mathbb{F}_q^{k'}$, with codeword
    $\mathrm{Encode}_{s,d,k',q}(F) = (J_{\mathbf{x}}^{<s}(F))_{\mathbf{x} \in
        \mathbb{F}_q^{k'}}$ for each $F$ of total degree at most $d$.
\end{definition}

The key analytic fact is Hermite interpolation: over any field, a univariate
polynomial of degree at most $d$ is determined by its order-$<s$ Hasse jets at
enough points.

\begin{lemma}[Hermite interpolation of Hasse
    derivatives~{\cite[Lem.~2.8]{GIJS26}}]
    Let $f$ be a univariate polynomial of degree at most $d$ over a finite
    field
    $\mathbb{F}$, let $\alpha_1,\dots,\alpha_h$ be distinct elements of
    $\mathbb{F}$, and let $\mu_1,\dots,\mu_h$ be positive integers with
    $\sum_i \mu_i > d$.  Then $f$ is recovered in
    $\mathrm{poly}(d,h,\log |\mathbb{F}|)$ time from the values
    $(\partial^{[j]} f)(\alpha_i)$ for $i \in [h]$, $0 \leq j < \mu_i$.
\end{lemma}

\paragraph{Use in the GIJS distinguisher.}
Let $Y = [\mathbf{y}_1 \mid \cdots \mid \mathbf{y}_{n'}] \in
    \mathbb{F}_q^{k' \times n'}$ be the input matrix and fix a hold-out column
$\tau$.  The distinguisher computes the kernel
\[
    K_\tau := \Bigl\{H \in \mathcal{H}_{k',d} :
    J_{\mathbf{y}_j}^{<s}(H) = 0 \text{ for all } j \neq \tau\Bigr\},
\]
where $\mathcal{H}_{k',d}$ is the space of \emph{homogeneous} degree-$d$
polynomials in $k'$ variables, of dimension $\binom{k'+d-1}{d}$.  It outputs
``GRS'' iff every $H \in K_\tau$ vanishes at $\mathbf{y}_\tau$.  Restricting
to homogeneous polynomials kills the unknown GRS column multipliers: for
homogeneous $H$ of degree $d$,
\[
    (\partial^{[\mathbf{a}]} H)(\lambda \mathbf{x})
    = \lambda^{d - |\mathbf{a}|} \cdot (\partial^{[\mathbf{a}]} H)(\mathbf{x})
    \qquad (\lambda \neq 0),
\]
so $J_{\lambda \mathbf{x}}^{<s}(H) = 0 \iff J_{\mathbf{x}}^{<s}(H) = 0$
\cite[Lem.~3.1]{GIJS26}.  The two feasibility conditions are the
\emph{reconstruction} inequality $s(n'-1) > dD$, which forces every
$H \in K_\tau$ to vanish at a clean hold-out via Hermite interpolation along
the hidden degree-$D$ curve, and the \emph{kernel} inequality
$(n'-1)\binom{k'+s-1}{s-1} < \binom{k'+d-1}{d}$, which guarantees
$K_\tau \neq \{0\}$ on a random input so that Schwartz--Zippel applies.


\subsection{The GIJS distinguisher}\label{sec:gijs-dist}


The Ghoshal--Ishai--Jain--Sun~\cite{GIJS26} distinguisher tests whether a given matrix spans
a subcode of a generalized Reed--Solomon code. Given $n'$ columns in
$\mathbb{F}_q^{k'}$, it interpolates: it seeks a nonzero homogeneous degree-$d$
polynomial vanishing to order $s$ at $n'-1$ of the columns (the interpolation
set), then checks whether it also vanishes at the remaining hold-out column.
For a true GRS subcode, the hidden rational-curve structure forces this
vanishing; for random columns, Schwartz--Zippel says it fails with high
probability. The two inequalities below encode feasibility: the first ensures
enough interpolation constraints to force vanishing on the curve, the second
that the interpolation kernel is nonzero.

\begin{theorem}[Ghoshal et al., Theorem~3.2]\label{thm:ghoshal}
    Let $n' \leq q$, $k' \leq K$, $d < q$, $D := K-1$. If
    \begin{align}
        s(n'-1)                   & > d\cdot D \label{eq:ghoshal1}          \\
        (n'-1)\binom{k'+s-1}{s-1} & < \binom{k'+d-1}{d} \label{eq:ghoshal2}
    \end{align}
    then there is a $\mathrm{poly}\!\left(\binom{k'+d-1}{d}, n', \log
        q\right)$-time
    algorithm distinguishing a random $k'$-dimensional subcode of random
    $GRS_{n',K}$
    from uniform random, with gap $\geq 1 - d/q$.
\end{theorem}
The algorithm (Algorithm~1 of~\cite{GIJS26}): hold out an arbitrary column
$\tau$, compute
the kernel $K_\tau$ of Hasse-jet evaluations (order $< s$) at the other $n'-1$
columns in the space of homogeneous degree-$d$ polynomials, and output ``GRS''
iff every $H \in K_\tau$ vanishes at $y_\tau$, ``Random'' otherwise.

\paragraph{Variables.} $k'$ is the dimension of the subcode; $K$ the
hidden ambient GRS dimension ($k' + w$ here); $s$ the jet multiplicity; $d$ the
test-polynomial degree. The first inequality is \emph{reconstruction}; the
second the \emph{kernel} condition.